% ============================================================
%  Guía de aplicación: un tablero para la gerencia de Cafés Altavista
%  Semana 8 · Clase 16: MasterClass (Recapitulación del Módulo 3)
%  Curso: Análisis y Visualización de Datos para la Toma de Decisiones
%  Datos: datos/cafes_altavista.xlsx (simulados; empresa ficticia)
%  Compilar con: xelatex guia-aplicacion-modulo3.tex (dos veces)
% ============================================================
\documentclass[11pt]{article}

\usepackage[a4paper,margin=2cm]{geometry}
\usepackage{fontspec}
\usepackage{amssymb}
\usepackage[spanish,es-noshorthands,es-tabla]{babel}
\usepackage{xcolor}
\usepackage{graphicx}
\usepackage{booktabs}
\usepackage{array}
\usepackage{tabularx}
\usepackage{enumitem}
\usepackage{needspace}
\usepackage{tcolorbox}
\usepackage{tikz}
\usepackage{fancyhdr}
\usepackage{hyperref}

\tcbuselibrary{skins,breakable}
\usetikzlibrary{arrows.meta,positioning,calc,shapes.geometric,shapes.multipart,babel}

% ---------------- Tipografía (misma lógica que la plantilla de slides) ----------------
\IfFontExistsTF{Aptos}{\setmainfont{Aptos}}{\setmainfont{Lato}}
\setmonofont{DejaVu Sans Mono}[Scale=0.82]

% ---------------- Paleta Esumer ----------------
\definecolor{esumerNavy}{HTML}{262E66}
\definecolor{esumerNavyDark}{HTML}{1E4093}
\definecolor{esumerTeal}{HTML}{0CB3B6}
\definecolor{esumerSky}{HTML}{26A1DB}
\definecolor{esumerOrange}{HTML}{EF7C2C}
\definecolor{esumerGreen}{HTML}{94C122}
\definecolor{esumerMagenta}{HTML}{AF348B}
\definecolor{esumerGray}{HTML}{E8E8E8}
\definecolor{esumerBg}{HTML}{F7F8FA}

\hypersetup{colorlinks=true, linkcolor=esumerNavyDark, urlcolor=esumerNavyDark}
\graphicspath{{../../../plantilla/assets/}}

% ---------------- Enlaces ----------------
\newcommand{\urldatos}{https://github.com/craljimenez/An-lisis_y_Visualizacion_de_Datos_para_la_Toma_de_Decisiones/raw/main/clases/semana08-clase16-masterclass-recapitulacion-modulo3/datos/cafes_altavista.xlsx}
\newcommand{\urlstore}{https://apps.microsoft.com/detail/9ntxr16hnw1t?hl=es-co}
\newcommand{\urlconsultas}{https://learn.microsoft.com/es-es/power-bi/transform-model/desktop-query-overview}
\newcommand{\urlrelaciones}{https://learn.microsoft.com/es-es/power-bi/transform-model/desktop-create-and-manage-relationships}
\newcommand{\urlestrella}{https://learn.microsoft.com/es-es/power-bi/guidance/star-schema}
\newcommand{\urlcolumnas}{https://learn.microsoft.com/es-es/power-bi/transform-model/desktop-calculated-columns}
\newcommand{\urlmedidas}{https://learn.microsoft.com/es-es/power-bi/transform-model/desktop-measures}
\newcommand{\urlkpi}{https://learn.microsoft.com/es-es/power-bi/visuals/power-bi-visualization-kpi}
\newcommand{\urlsegmentadores}{https://learn.microsoft.com/es-es/power-bi/visuals/power-bi-visualization-slicers}
\newcommand{\urlfiltros}{https://learn.microsoft.com/es-es/power-bi/create-reports/power-bi-report-add-filter}
\newcommand{\urlinteracciones}{https://learn.microsoft.com/es-es/power-bi/create-reports/service-reports-visual-interactions}
\newcommand{\urlexplorar}{https://learn.microsoft.com/es-es/power-bi/explore-reports/end-user-drill}
\newcommand{\urlazuremaps}{https://learn.microsoft.com/es-es/azure/azure-maps/power-bi-visual-get-started}
\newcommand{\urltrucos}{https://learn.microsoft.com/es-es/power-bi/create-reports/desktop-tips-and-tricks-for-creating-reports}

% ---------------- Encabezado / pie de página ----------------
\pagestyle{fancy}
\fancyhf{}
\renewcommand{\headrulewidth}{0pt}
\fancyfoot[L]{\small\color{esumerNavyDark} Guía de aplicación del Módulo 3 \textperiodcentered\ Clase 16}
\fancyfoot[R]{\small\color{esumerNavyDark} Cristian A. Jiménez C. \quad\textperiodcentered\quad \thepage}
\setlength{\parindent}{0pt}
\setlength{\parskip}{4pt}

% ---------------- Fórmulas DAX ----------------
\newtcolorbox{dax}[1][]{
  colback=esumerBg, colframe=esumerNavy!35, boxrule=0.5pt, arc=1mm,
  left=6pt, right=6pt, top=3pt, bottom=3pt, fonttitle=\bfseries\small,
  coltitle=esumerNavy, colbacktitle=esumerGray, fontupper=\small\ttfamily, #1
}

% ---------------- Cajas reutilizables ----------------
\newtcolorbox{recuerda}[1][Lo viste en clase]{
  colback=esumerBg, colframe=esumerTeal, boxrule=0.6pt, arc=1mm,
  fonttitle=\bfseries\small, coltitle=white, colbacktitle=esumerTeal,
  title=#1, left=6pt, right=6pt, top=3pt, bottom=3pt
}
\newtcolorbox{cuidado}[1][Cuidado]{
  colback=white, colframe=esumerOrange, boxrule=0.6pt, arc=1mm,
  fonttitle=\bfseries\small, coltitle=white, colbacktitle=esumerOrange,
  title=#1, left=6pt, right=6pt, top=3pt, bottom=3pt
}
\newtcolorbox{control}[1][Cifras de control: compara antes de seguir]{
  colback=white, colframe=esumerGreen!80!black, boxrule=0.6pt, arc=1mm,
  fonttitle=\bfseries\small, coltitle=white, colbacktitle=esumerGreen!80!black,
  title=#1, left=6pt, right=6pt, top=3pt, bottom=3pt
}

% Menú de Power BI: nombre en español (nombre en inglés)
\newcommand{\menu}[2]{\textbf{#1} {\color{black!55}(#2)}}
% Referencia a la clase donde se vio
\newcommand{\vez}[1]{\leavevmode\mbox{\color{esumerTeal}\textbf{[#1]}}}

% ---------------- Encabezado de cada etapa ----------------
\newcommand{\etapa}[3]{%
  \clearpage\noindent
  \begin{tikzpicture}[baseline=(n.base)]
    \node[circle, fill=esumerOrange, text=white, font=\bfseries\large,
          minimum size=0.85cm, inner sep=0pt] (n) {#1};
  \end{tikzpicture}\hspace{8pt}%
  {\Large\bfseries\color{esumerNavy} #2}\par
  \vspace{-2pt}{\color{esumerOrange}\rule{\textwidth}{1pt}}\par
  \vspace{2pt}
  {\color{esumerNavyDark}\textbf{Objetivo:}} \textit{#3}\par
  \vspace{2pt}
}
\newcommand{\subtitulo}[1]{\par\needspace{3\baselineskip}\vspace{4pt}{\large\bfseries\color{esumerNavyDark} #1}\par}

% ---------------- Lista de chequeo ----------------
\newlist{checklist}{itemize}{1}
\setlist[checklist]{label=$\square$, itemsep=1pt, topsep=2pt, leftmargin=1.4em}

% ---------------- Estilos TikZ comunes ----------------
\tikzset{
  every node/.append style={execute at begin node={\hyphenpenalty=10000\exhyphenpenalty=10000\relax}},
  flecha/.style={-{Stealth[length=2.4mm]}, line width=0.9pt, color=esumerNavy},
  caja/.style={draw=esumerNavy, line width=0.7pt, rounded corners=3pt,
               fill=white, align=center, inner sep=4pt},
  num/.style={circle, fill=esumerOrange, text=white, font=\bfseries\small,
              minimum size=0.6cm, inner sep=0pt},
}

% ============================================================
\begin{document}

% ---------------- Encabezado ----------------
\begin{minipage}[c]{0.18\textwidth}
  \includegraphics[height=1.15cm,width=\linewidth,keepaspectratio]{esumer-logo.png}
\end{minipage}%
\begin{minipage}[c]{0.80\textwidth}
  \raggedleft
  {\Large\bfseries\color{esumerNavy} Guía de aplicación: un tablero para la gerencia}\\[3pt]
  {\small\color{esumerNavyDark} Semana 8 \textperiodcentered\ Clase 16 (MasterClass del Módulo 3) \textperiodcentered\ Análisis y Visualización de Datos para la Toma de Decisiones}
\end{minipage}

\vspace{4pt}
{\color{esumerOrange}\rule{\textwidth}{2pt}}
\vspace{6pt}

\textbf{Para qué es esta guía.} En las Clases 13, 14 y 15 comenzamos a construir, paso a
paso y con instrucciones detalladas, un tablero del gimnasio. Hoy haces lo
mismo \textbf{por tu cuenta}, con datos que no conoces: los de una cadena de
cafeterías. La guía te dice qué construir y cómo comprobar que vas bien,
pero ya no te da cada clic: cuando dudes, las pistas entre corchetes,
como \vez{Clase 13, paso 5}, te dicen en qué clase y en qué paso lo viste.

\begin{recuerda}[El caso: Cafés Altavista S.A.S.]
\small Cafés Altavista es una cadena de \textbf{seis cafeterías} en Medellín,
Bogotá, Cali y Pereira. Vende bebidas calientes, bebidas frías, panadería y
café en grano para llevar a casa. Su gerencia te contrató para responder,
con un tablero de Power BI, las preguntas que llevará al comité de octubre
de 2026. Tienes un año de ventas, de octubre de 2025 a septiembre de 2026.

\vspace{2pt}
\footnotesize\itshape La empresa es ficticia y los datos son simulados con
fines académicos, pero tienen los problemas de un archivo real: errores de
digitación, duplicados, celdas vacías y alguna sorpresa.
\end{recuerda}

\textbf{Al terminar la sesión deberías tener:}
\begin{enumerate}[itemsep=1pt, topsep=2pt]
  \item Un archivo \texttt{cafes\_altavista.pbix} con los datos limpios en
    Power Query y un modelo estrella con relaciones correctas.
  \item Las medidas DAX de la etapa 3, con las cifras de control comprobadas.
  \item Una página exploratoria y una \textbf{página para la gerencia} con
    KPI, segmentadores, una jerarquía y un mapa.
  \item Una \textbf{ficha de 3 a 5 insights} y una presentación de cinco
    minutos para el comité (etapa 6).
\end{enumerate}

\textbf{Agenda sugerida para las 3 horas}

\begin{tabular}{@{}>{\raggedright\arraybackslash}p{1.4cm}>{\raggedright\arraybackslash}p{6.9cm}>{\raggedright\arraybackslash}p{7.6cm}@{}}
\toprule
\textbf{Hora} & \textbf{Etapas} & \textbf{Producto al final de la hora}\\
\midrule
Hora 1 & 1. Cargar y limpiar \textperiodcentered\ 2. Modelo & Cuatro tablas limpias, \texttt{Calendario} y cuatro relaciones.\\
Hora 2 & 3. Columnas y medidas \textperiodcentered\ 4. Explorar & Medidas con cifras de control y respuestas a las preguntas.\\
Hora 3 & 5. Página para la gerencia \textperiodcentered\ 6. Comunicar & Página final, ficha de insights y presentación para el comité.\\
\bottomrule
\end{tabular}

\vspace{6pt}
\textbf{Antes de empezar:} descarga el archivo
\href{\urldatos}{\color{esumerOrange}\texttt{cafes\_altavista.xlsx}} (unos 4 MB)
y abre Power BI Desktop
(\href{\urlstore}{\color{esumerOrange}descarga en Microsoft Store} si no lo
tienes). Guarda tu archivo \texttt{.pbix} al terminar cada etapa.

% ============================================================
\clearpage
\section*{\color{esumerNavy}La carta de la gerencia}

Esto es lo que te escribió la gerente general. Léelo completo: cada
pregunta te dice qué columnas, qué medidas y qué gráfico vas a necesitar.

\begin{tcolorbox}[colback=white, colframe=esumerNavy, boxrule=0.8pt, arc=1mm,
  left=10pt, right=10pt, top=6pt, bottom=6pt]
\small
\textit{Hola. Para el comité de octubre necesito respuestas claras, con cifras,
a estas preguntas:}

\begin{enumerate}[label=\textbf{P\arabic*.}, leftmargin=2.2em, itemsep=3pt, topsep=4pt]
  \item \textit{En septiembre vendimos más que nunca. Pero el contador dice
    que el margen se ve peor. ¿Qué pasó?}
  \item \textit{¿Qué tiendas cumplen su meta mensual de ventas? Tengo la
    impresión de que Chapinero no va bien: ¿es cierto?, ¿cuándo le va peor?}
  \item \textit{Lanzamos la app el 1 de marzo de 2026. ¿Está funcionando? ¿Quienes
    piden por la app compran más o menos que en el mostrador?}
  \item \textit{¿A qué horas debemos reforzar el personal, y qué productos
    hay que tener listos en cada momento del día?}
  \item \textit{Usaquén abrió el 1 de abril de 2026. ¿Va por buen camino?}
  \item \textit{¿Vale la pena el programa de socios? ¿Los socios compran más
    que los clientes ocasionales?}
\end{enumerate}

\textit{Una cosa más: el reporte de julio de Granada, en Cali, me pareció
raro. Revísalo antes de presentarlo.}

\vspace{4pt}
\textit{Gracias.}\\
\textbf{Gerencia General, Cafés Altavista}
\end{tcolorbox}

\subtitulo{Lo que la gerencia ya sabe, y tú también}

\begin{itemize}[itemsep=2pt, topsep=2pt]
  \item \textbf{Metas de ventas.} Cada tienda tiene una meta \emph{mensual}
    de ventas netas. Está en la hoja \texttt{Tiendas}, en la columna
    \texttt{meta\_ventas\_mensual\_cop}.
  \item \textbf{Meta de margen.} La empresa espera un margen bruto de al
    menos el \textbf{60\%} de las ventas netas. Es una decisión de la
    gerencia, no sale de los datos (Clase 15, paso 4).
  \item \textbf{Campaña de septiembre.} En septiembre de 2026 se hizo la campaña
    ``Mes del grano'': 25\% de descuento en café en grano.
  \item \textbf{Socios.} Los clientes inscritos en el programa de socios
    tienen un \texttt{id\_cliente}. Si la venta no tiene \texttt{id\_cliente},
    fue a un cliente ocasional. Los socios reciben descuentos en algunas
    compras.
  \item \textbf{Horario.} Todas las tiendas abren de 6:30 a.\,m. a 8:00 p.\,m.
\end{itemize}

\subtitulo{Cómo trabajar}

\begin{itemize}[itemsep=2pt, topsep=2pt]
  \item \textbf{En parejas}, si el grupo lo permite: uno construye y el otro
    revisa las cifras de control. Cambien de rol en cada etapa.
  \item \textbf{No avances si una cifra de control no cuadra.} Casi siempre
    el error está en la etapa anterior: un duplicado que no quitaste, una
    relación que falta o un tipo de dato equivocado. La última página tiene
    una tabla con los errores más comunes.
  \item \textbf{Cada pregunta de la gerencia termina en una frase}, no en un
    gráfico: una cifra, lo que significa y qué recomiendas hacer.
\end{itemize}

% ============================================================
\clearpage
\section*{\color{esumerNavy}Los datos: \texttt{cafes\_altavista.xlsx}}

Un libro de Excel con cuatro hojas, tal como las exporta el sistema de caja.
\textbf{No las corrijas en Excel}: toda la limpieza se hace en Power Query,
para que quede registrada en los pasos aplicados (Clase 13).

\vspace{2pt}
\small
\renewcommand{\arraystretch}{1.12}
\begin{tabularx}{\textwidth}{@{}>{\raggedright\arraybackslash}p{4.4cm} >{\raggedright\arraybackslash}X@{}}
\toprule
\textbf{Columna} & \textbf{Qué contiene}\\
\midrule
\multicolumn{2}{@{}l}{\color{esumerOrange}\textbf{Ventas}: una fila por cada producto de cada compra (77.860 filas en el Excel)}\\
\texttt{id\_linea} & Identificador de la fila. No debería repetirse.\\
\texttt{id\_ticket} & Identificador de la compra. Un ticket puede tener varias líneas (un capuchino y un croissant).\\
\texttt{fecha}, \texttt{hora} & Día y hora de la compra.\\
\texttt{id\_tienda}, \texttt{id\_producto} & Tienda donde se vendió y producto vendido.\\
\texttt{id\_cliente} & Socio que compró. Vacío si fue un cliente ocasional.\\
\texttt{canal} & \texttt{Mostrador}, \texttt{Domicilio} o \texttt{App}.\\
\texttt{cantidad}, \texttt{precio\_unitario\_cop} & Unidades y precio de cada una, en pesos.\\
\texttt{descuento\_cop} & Descuento de la línea, en pesos. Vacío si no hubo descuento.\\
\texttt{medio\_pago} & Efectivo, tarjeta, Nequi\ldots\\
\midrule
\multicolumn{2}{@{}l}{\color{esumerOrange}\textbf{Productos}: una fila por producto (18 productos)}\\
\texttt{id\_producto}, \texttt{producto} & Código y nombre.\\
\texttt{categoria}, \texttt{subcategoria} & Cuatro categorías, cada una con sus subcategorías (por ejemplo, Bebidas calientes $>$ Café con leche).\\
\texttt{precio\_lista\_cop} & Precio de venta normal, por unidad.\\
\texttt{costo\_unitario\_cop} & Lo que le cuesta a la empresa cada unidad (insumos y empaque).\\
\midrule
\multicolumn{2}{@{}l}{\color{esumerOrange}\textbf{Tiendas}: una fila por tienda (6 tiendas)}\\
\texttt{id\_tienda}, \texttt{tienda} & Código y nombre.\\
\texttt{ciudad}, \texttt{departamento}, \texttt{pais} & Ubicación por nombre.\\
\texttt{latitud}, \texttt{longitud} & Ubicación por coordenadas.\\
\texttt{fecha\_apertura}, \texttt{area\_m2} & Desde cuándo funciona y su tamaño.\\
\texttt{meta\_ventas\_mensual\_cop} & Meta de ventas netas de \emph{un mes}, en pesos.\\
\midrule
\multicolumn{2}{@{}l}{\color{esumerOrange}\textbf{Clientes}: una fila por socio del programa (905 filas en el Excel)}\\
\texttt{id\_cliente}, \texttt{nombre} & Código y nombre del socio.\\
\texttt{fecha\_registro}, \texttt{fecha\_nacimiento} & Cuándo se inscribió y cuándo nació.\\
\texttt{ciudad}, \texttt{correo} & Ciudad y correo que dio al inscribirse.\\
\texttt{acepta\_promociones} & Si autorizó recibir promociones.\\
\bottomrule
\end{tabularx}
\renewcommand{\arraystretch}{1.0}
\normalsize

\vspace{4pt}
\textbf{Dos definiciones que vas a usar todo el tiempo:}

\begin{center}
\begin{tikzpicture}[
  term/.style={caja, minimum height=1.1cm, font=\small},
  op/.style={font=\large\bfseries, text=esumerNavy},
]
  \node[term, fill=esumerTeal!15, text width=2.6cm] (vn) at (1.5,0) {\textbf{Venta neta}\\de una línea};
  \node[op] at (3.25,0) {=};
  \node[term, text width=1.6cm] at (4.6,0) {cantidad};
  \node[op] at (5.75,0) {$\times$};
  \node[term, text width=2.7cm] at (7.6,0) {precio unitario};
  \node[op] at (9.5,0) {$-$};
  \node[term, text width=2.2cm] at (11.1,0) {descuento};
  \node[term, fill=esumerTeal!15, text width=2.6cm] (mg) at (1.5,-1.5) {\textbf{Margen}};
  \node[op] at (3.25,-1.5) {=};
  \node[term, text width=2.5cm] at (5.0,-1.5) {venta neta};
  \node[op] at (6.75,-1.5) {$-$};
  \node[term, text width=3.9cm] at (9.3,-1.5) {cantidad $\times$ costo unitario};
  \node[font=\footnotesize, text width=3.3cm, align=left] at (14.8,-0.75)
    {El \textbf{margen \%} es el margen dividido por la venta neta: de cada
     peso vendido, cuánto queda después de pagar el producto.};
\end{tikzpicture}
\end{center}

% ============================================================
\clearpage
\section*{\color{esumerNavy}El recorrido completo}

Seis etapas, las mismas del Módulo 3. Las cuatro primeras construyen; las
dos últimas comunican. La etiqueta de cada caja te dice en qué clase lo
viste.

\vspace{4pt}
\begin{center}
\begin{tikzpicture}[
  et/.style={draw=esumerNavy, line width=0.8pt, rounded corners=3pt,
             fill=white, minimum width=5.0cm, minimum height=3.1cm,
             text width=4.7cm, align=center, inner sep=3pt, font=\scriptsize},
  clase/.style={fill=esumerTeal, text=white, font=\scriptsize\bfseries,
                rounded corners=2pt, inner sep=2.5pt},
]
  \node[et] (e1) at (2.85,0)
    {{\bfseries\small Cargar y limpiar}\\[3pt]
     {Power Query: tipos, duplicados, vacíos, texto con
      variantes\\[2pt]\itshape ¿Mis datos dicen lo que creo que dicen?}};
  \node[et] (e2) at (8.5,0)
    {{\bfseries\small Modelo de datos}\\[3pt]
     {Hechos y dimensiones, \texttt{Calendario}, relaciones
      muchos a uno\\[2pt]\itshape ¿Cada filtro llega a donde debe?}};
  \node[et] (e3) at (14.15,0)
    {{\bfseries\small Columnas y medidas}\\[3pt]
     {DAX: venta neta, costo, margen, ticket, meta,
      variación\\[2pt]\itshape ¿Mis cifras cuadran con las de control?}};
  \node[et] (e4) at (14.15,-4.0)
    {{\bfseries\small Explorar}\\[3pt]
     {Barras, líneas, columnas, mapa, anomalías\\[2pt]
      \itshape ¿Qué responde cada pregunta de la gerencia?}};
  \node[et] (e5) at (8.5,-4.0)
    {{\bfseries\small Página para la gerencia}\\[3pt]
     {KPI frente a la meta, segmentadores, jerarquía,
      interacciones\\[2pt]\itshape ¿Otro la entiende sin que yo la explique?}};
  \node[et] (e6) at (2.85,-4.0)
    {{\bfseries\small Comunicar}\\[3pt]
     {Ficha de insights y presentación de cinco
      minutos\\[2pt]\itshape ¿Qué decisión se toma con esto?}};
  \foreach \i in {1,...,6} {\node[num] at (e\i.north west) {\i};}
  \node[clase, anchor=south east] at ($(e1.north east)+(-0.15,-0.02)$) {Clase 13};
  \node[clase, anchor=south east] at ($(e2.north east)+(-0.15,-0.02)$) {Clase 13};
  \node[clase, anchor=south east] at ($(e3.north east)+(-0.15,-0.02)$) {Clases 13 a 15};
  \node[clase, anchor=south east] at ($(e4.north east)+(-0.15,-0.02)$) {Clase 14};
  \node[clase, anchor=south east] at ($(e5.north east)+(-0.15,-0.02)$) {Clase 15};
  \node[clase, anchor=south east] at ($(e6.north east)+(-0.15,-0.02)$) {Clases 11 y 14};
  \draw[flecha] (e1) -- (e2);
  \draw[flecha] (e2) -- (e3);
  \draw[flecha] (e3) -- (e4);
  \draw[flecha] (e4) -- (e5);
  \draw[flecha] (e5) -- (e6);
  \node[draw=esumerGreen!80!black, dashed, line width=0.8pt, rounded corners=3pt,
        fill=esumerBg, text width=16.1cm, align=left, inner sep=5pt]
    at (8.5,-6.75)
    {\small\textbf{\color{esumerNavy}Cifras de control.} Al final de cada
     etapa hay una caja verde con números que tu archivo debe mostrar. Si
     coinciden, sigue. Si no, busca la causa antes de avanzar: un error en
     la etapa 1 se arrastra hasta la presentación.};
\end{tikzpicture}
\end{center}

\vspace{6pt}
Tres reglas para todo el recorrido:

\begin{enumerate}[itemsep=2pt, topsep=2pt]
  \item \textbf{Nunca edites el Excel.} Si te equivocas en Power Query,
    borra el paso con la \textbf{X} en \menu{Pasos aplicados}{Applied steps}
    y vuelve a hacerlo. El archivo original queda intacto.
  \item \textbf{Una medida, no una columna, para lo que se suma o se
    divide.} El margen \% de un mes no es el promedio de los márgenes de cada
    fila: es el margen del mes dividido por las ventas del mes. Eso lo hace
    bien una medida (Clase 13, paso 18).
  \item \textbf{Antes de concluir, valida.} Una cifra rara puede ser un
    hallazgo o un error de datos. Pregunta: ¿pasa en todas las tiendas?, ¿en
    qué días exactos?, ¿cuadra con otra fuente? (Clase 14).
\end{enumerate}

% ============================================================
\etapa{1}{Cargar y limpiar en Power Query}{dejar las cuatro tablas limpias, con los tipos correctos y sin duplicados ni variantes de texto.}

\textbf{Cargar.} \menu{Obtener datos}{Get data} $>$ \menu{Libro de
Excel}{Excel workbook}, elige el archivo y marca las \textbf{cuatro hojas} en
el Navegador. Pulsa \menu{Transformar datos}{Transform data}, no
\menu{Cargar}{Load}: primero se limpia. \vez{Clase 13, pasos 1 a 3}

\textbf{Diagnosticar.} En la pestaña \menu{Vista}{View} activa
\menu{Calidad de columnas}{Column quality}, \menu{Distribución de
columnas}{Column distribution} y \menu{Perfil de columna}{Column profile}.
Abajo a la izquierda, cambia a \textbf{basado en todo el conjunto de datos}:
si no, solo ves las primeras 1000 filas. \vez{Clase 13, paso 4}

\subtitulo{Lo que tienes que hacer en cada hoja}

\small
\renewcommand{\arraystretch}{1.2}
\begin{tabularx}{\textwidth}{@{}>{\raggedright\arraybackslash}p{2.1cm} >{\raggedright\arraybackslash}X >{\raggedright\arraybackslash}p{3.5cm}@{}}
\toprule
\textbf{Hoja} & \textbf{Tarea} & \textbf{Pista}\\
\midrule
\texttt{Ventas} & Quitar las filas repetidas: selecciona \texttt{id\_linea} $>$ \menu{Quitar filas}{Remove rows} $>$ \menu{Quitar duplicados}{Remove duplicates}. & \vez{Clase 13, paso 5}\\
& Tipos: \texttt{fecha} como \textbf{Fecha} (no Fecha y hora), \texttt{hora} como \textbf{Hora}, los \texttt{id\_} y los valores en pesos como \textbf{Número entero}. & \vez{Clase 13, paso 6}\\
& En \texttt{descuento\_cop}, reemplazar \texttt{null} por \texttt{0}: un vacío aquí significa ``sin descuento''. & \vez{Clase 13, paso 9}\\
\midrule
\texttt{Productos} & Unificar \texttt{categoria}: tiene 11 formas de escribir 4 categorías (ver el diagrama). & \vez{Clase 13, pasos 7 a 9}\\
\midrule
\texttt{Clientes} & Quitar los socios repetidos (por \texttt{id\_cliente}). & \vez{Clase 13, paso 5}\\
& Unificar \texttt{ciudad}: \texttt{Bogota}, \texttt{Bogotá D.C.} y \texttt{BOGOTÁ} son la misma ciudad; \texttt{Santiago de Cali} es Cali. & \vez{Clase 13, pasos 7 a 9}\\
& Unificar \texttt{acepta\_promociones}: hay \texttt{Sí}, \texttt{si}, \texttt{SI}, \texttt{1}, \texttt{No}, \texttt{no}, \texttt{NO}, \texttt{0}. Cámbiala primero a tipo \textbf{Texto}. & \vez{Clase 13, pasos 6 a 9}\\
& \texttt{correo} y \texttt{fecha\_nacimiento} tienen vacíos: \textbf{déjalos}. No hay un valor razonable para inventar un correo, y no los usarás en el tablero. & \vez{Clase 8}\\
\midrule
\texttt{Tiendas} & Revisar tipos: \texttt{fecha\_apertura} como Fecha; \texttt{latitud} y \texttt{longitud} como \textbf{Número decimal}. & \vez{Clase 13, paso 6}\\
\bottomrule
\end{tabularx}
\renewcommand{\arraystretch}{1.0}
\normalsize

\subtitulo{Cómo unificar un texto con variantes: el mismo método de \texttt{sexo}}

\begin{center}
\begin{tikzpicture}[
  col/.style={draw=esumerNavy!50, fill=white, rounded corners=2pt, text width=3.4cm,
              align=center, font=\scriptsize\ttfamily, inner sep=4pt},
  tit/.style={font=\bfseries\scriptsize, text=esumerNavy, align=center, text width=3.8cm},
  fl/.style={-{Stealth[length=2.2mm]}, line width=0.9pt, color=esumerOrange},
]
  \node[col] (a) at (1.95,0) {"Bebidas calientes"\\"Bebidas Calientes "\\"bebidas calientes"\\[2pt]
                             "Bebidas frías" "Bebidas frias"\\"BEBIDAS FRÍAS"\\[2pt]
                             "Café en grano"\\"Cafe en grano"\\[2pt]
                             "Panadería" "Panaderia"\\" panadería"};
  \node[col] (b) at (6.35,0) {"BEBIDAS CALIENTES"\\[2pt]
                             "BEBIDAS FRÍAS"\\"BEBIDAS FRIAS"\\[2pt]
                             "CAFÉ EN GRANO"\\"CAFE EN GRANO"\\[2pt]
                             "PANADERÍA"\\"PANADERIA"};
  \node[col] (c) at (10.75,0) {"BEBIDAS CALIENTES"\\[2pt]"BEBIDAS FRÍAS"\\[2pt]
                              "CAFÉ EN GRANO"\\[2pt]"PANADERÍA"};
  \node[col] (d) at (15.05,0) {"Bebidas Calientes"\\[2pt]"Bebidas Frías"\\[2pt]
                              "Café En Grano"\\[2pt]"Panadería"};
  \node[tit, above=0.15cm of a] {En el Excel (11 variantes)};
  \node[tit, above=0.15cm of b] {Tras Recortar y MAYÚSCULAS};
  \node[tit, above=0.15cm of c] {Tras Reemplazar valores};
  \node[tit, text width=3.0cm, above=0.15cm of d] {Tras Poner En Mayúsculas Cada Palabra};
  \draw[fl] (a) -- (b);
  \draw[fl] (b) -- (c);
  \draw[fl] (c) -- (d);
\end{tikzpicture}
\end{center}

\small Los tres reemplazos: \texttt{BEBIDAS FRIAS} por \texttt{BEBIDAS FRÍAS},
\texttt{PANADERIA} por \texttt{PANADERÍA} y \texttt{CAFE EN GRANO} por
\texttt{CAFÉ EN GRANO}. Haz lo mismo con \texttt{ciudad}
(\texttt{BOGOTA}, \texttt{BOGOTÁ D.C.}, \texttt{MEDELLIN} y
\texttt{SANTIAGO DE CALI}) y con \texttt{acepta\_promociones}
(\texttt{SI}, \texttt{1} y \texttt{0}).
\normalsize

\begin{cuidado}[Marca ``Coincidir con todo el contenido de la celda'']
\small En \menu{Opciones avanzadas}{Advanced options} de Reemplazar valores.
Sin esa casilla, Power Query cambia el texto aunque esté dentro de otro:
al reemplazar \texttt{1} por \texttt{SÍ}, un \texttt{10} quedaría
\texttt{SÍ0}. \vez{Clase 13, paso 9}
\end{cuidado}

\subtitulo{Cerrar y comprobar}

Cuando las cuatro tablas estén listas: \menu{Inicio}{Home} $>$
\menu{Cerrar y aplicar}{Close \& apply}. La carga de \texttt{Ventas} tarda
unos segundos. Después, en la vista \textbf{Tabla} (la segunda de la barra
izquierda), elige cada tabla en el panel \menu{Datos}{Data}: abajo a la
izquierda aparece cuántas filas tiene. \vez{Clase 13, paso 13}

\begin{control}
\small
\renewcommand{\arraystretch}{1.15}
\begin{tabularx}{\linewidth}{@{}>{\raggedright\arraybackslash}p{3.4cm} >{\raggedright\arraybackslash}p{3.1cm} >{\raggedright\arraybackslash}X@{}}
\toprule
\textbf{Tabla} & \textbf{Filas} & \textbf{Qué cambió}\\
\midrule
\texttt{Ventas} & \textbf{77.840} & Se quitaron 20 filas repetidas (las tenía el Excel: 77.860).\\
\texttt{Productos} & \textbf{18} & \texttt{categoria}: 6 Bebidas Calientes, 4 Bebidas Frías, 3 Café En Grano y 5 Panadería.\\
\texttt{Tiendas} & \textbf{6} & Sin cambios en las filas.\\
\texttt{Clientes} & \textbf{900} & Se quitaron 5 socios repetidos. \texttt{ciudad}: Medellín 356, Bogotá 326, Cali 125 y Pereira 93. \texttt{acepta\_promociones}: Sí 536 y No 364.\\
\bottomrule
\end{tabularx}
\end{control}

\vspace{2pt}
Para ver los conteos de una columna, en Power Query activa el
\menu{Perfil de columna}{Column profile} y selecciónala: la
\menu{Distribución de valores}{Value distribution} muestra cuántas filas
tiene cada valor.

\begin{recuerda}[¿De dónde salieron las 20 filas repetidas?]
\small No borres sin preguntar: identifica qué filas eran. En el Excel crudo,
filtra \texttt{id\_linea} por un valor repetido y mira la \texttt{fecha} y
la tienda. Si todas las repetidas son del mismo día y de la misma tienda, lo
más probable es que ese día se exportó dos veces. Anótalo: es un hallazgo
sobre la \emph{calidad} de los datos y va en tu ficha (etapa 6).
\end{recuerda}

\begin{checklist}
  \item Las cuatro tablas tienen las filas de la caja verde.
  \item En \texttt{Ventas}, \texttt{fecha} tiene un ícono de calendario y
    \texttt{hora} un ícono de reloj.
  \item En \texttt{descuento\_cop} ya no hay \texttt{null}.
  \item \texttt{categoria}, \texttt{ciudad} y \texttt{acepta\_promociones}
    tienen 4, 4 y 2 valores distintos.
  \item Guardé el archivo como \texttt{cafes\_altavista.pbix}.
\end{checklist}

% ============================================================
\etapa{2}{El modelo de datos}{armar un modelo estrella para que cada filtro (tienda, producto, socio o fecha) llegue a la tabla de ventas.}

\texttt{Ventas} es la tabla de \textbf{hechos}: cada fila es algo que pasó.
\texttt{Productos}, \texttt{Tiendas} y \texttt{Clientes} son
\textbf{dimensiones}: describen el quién, el qué y el dónde. Falta el
\textbf{cuándo}: una tabla \texttt{Calendario} que creas con DAX.
\vez{Clase 13, ``Hechos y dimensiones''}

\begin{center}
\begin{tikzpicture}[
  dim/.style={rectangle split, rectangle split parts=2, draw=esumerNavy, line width=0.7pt,
              rounded corners=2pt, rectangle split part fill={esumerNavy,white},
              text width=3.7cm, align=left, font=\scriptsize\ttfamily, inner sep=3pt},
  hecho/.style={rectangle split, rectangle split parts=2, draw=esumerOrange, line width=0.9pt,
              rounded corners=2pt, rectangle split part fill={esumerOrange,white},
              text width=3.9cm, align=left, font=\scriptsize\ttfamily, inner sep=3pt},
  card/.style={font=\scriptsize\bfseries, text=esumerOrange},
]
  \node[hecho] (v) at (8.5,0) {\textcolor{white}{\bfseries\small\sffamily Ventas (hechos)}
    \nodepart{two}id\_linea, id\_ticket\\fecha, hora\\id\_tienda\\id\_producto\\id\_cliente\\canal, cantidad\\precio\_unitario\_cop\\descuento\_cop, medio\_pago};
  \node[dim] (p) at (2.3,1.4) {\textcolor{white}{\bfseries\small\sffamily Productos}
    \nodepart{two}id\_producto\\producto\\categoria, subcategoria\\costo\_unitario\_cop};
  \node[dim] (t) at (14.7,1.4) {\textcolor{white}{\bfseries\small\sffamily Tiendas}
    \nodepart{two}id\_tienda\\tienda, ciudad\\latitud, longitud\\meta\_ventas\_mensual\_cop};
  \node[dim] (c) at (2.3,-1.9) {\textcolor{white}{\bfseries\small\sffamily Clientes}
    \nodepart{two}id\_cliente\\nombre, ciudad\\acepta\_promociones};
  \node[dim, rectangle split part fill={esumerTeal,white}, draw=esumerTeal] (k) at (14.7,-1.9)
    {\textcolor{white}{\bfseries\small\sffamily Calendario (DAX)}
    \nodepart{two}Date\\Año, Mes, Año-Mes\\Inicio de mes, Día semana};
  \draw[flecha] (p.east) -- node[pos=0.12, above, card]{1} node[pos=0.88, above, card]{*} ($(v.north west)!0.45!(v.west)$);
  \draw[flecha] (t.west) -- node[pos=0.12, above, card]{1} node[pos=0.88, above, card]{*} ($(v.north east)!0.45!(v.east)$);
  \draw[flecha] (c.east) -- node[pos=0.12, above, card]{1} node[pos=0.88, above, card]{*} ($(v.south west)!0.45!(v.west)$);
  \draw[flecha] (k.west) -- node[pos=0.12, above, card]{1} node[pos=0.88, above, card]{*} ($(v.south east)!0.45!(v.east)$);
\end{tikzpicture}

{\footnotesize\itshape Las flechas muestran la dirección del filtro: de la
dimensión (lado 1) a los hechos (lado *).}
\end{center}

\subtitulo{Paso a paso}

\begin{enumerate}[itemsep=3pt, topsep=2pt]
  \item \textbf{Calendario.} Vista Tabla $>$ \menu{Nueva tabla}{New table}:
\begin{dax}
Calendario = CALENDAR(DATE(2025, 10, 1), DATE(2026, 9, 30))
\end{dax}
    Márcala como tabla de fechas, con la columna \texttt{Date}. Luego agrégale
    cinco columnas (\menu{Nueva columna}{New column}):
\begin{dax}
Año = YEAR(Calendario[Date])\\
Mes = MONTH(Calendario[Date])\\
Año-Mes = FORMAT(Calendario[Date], "yyyy-mm")\\
Inicio de mes = DATE(YEAR(Calendario[Date]), MONTH(Calendario[Date]), 1)\\
Día semana = WEEKDAY(Calendario[Date], 2)
\end{dax}
    \footnotesize \texttt{WEEKDAY} con un \texttt{2} numera los días de
    \texttt{1} (lunes) a \texttt{7} (domingo). Es la única función nueva de
    hoy.\par\normalsize\hfill\vez{Clase 13, paso 16; Clase 14, paso 1}
  \item \textbf{Relaciones.} En la vista \textbf{Modelo}, es probable que Power BI
    ya haya creado tres (detecta los nombres \texttt{id\_} repetidos). Revisa que sean
    \textbf{muchos a uno (*:1)} con dirección de filtro \textbf{única}, y
    crea la cuarta: arrastra \texttt{Ventas[fecha]} sobre
    \texttt{Calendario[Date]}. \hfill \vez{Clase 13, pasos 14 y 15}
\end{enumerate}

\begin{cuidado}[Si la relación con \texttt{Calendario} no filtra]
\small Si \texttt{Ventas[fecha]} quedó como \textbf{Fecha y hora}, la
relación se crea pero puede no encontrar coincidencias. Vuelve a Power Query
y cámbiala a \textbf{Fecha}.
\end{cuidado}

\begin{control}
\small Inserta un objeto visual \menu{Tabla}{Table} con
\texttt{Tiendas[tienda]} y \texttt{Ventas[id\_linea]} (en el pozo de
valores, elige \menu{Recuento}{Count}). Las seis tiendas deben sumar
\textbf{77.840}, y \textbf{Altavista Usaquén} debe tener 4.987 líneas. Si
ves una sola cifra repetida en todas las filas, falta la relación
\texttt{Tiendas}--\texttt{Ventas}. \vez{Clase 13, paso 19}
\end{control}

% ============================================================
\etapa{3}{Columnas calculadas y medidas}{calcular, una sola vez y bien, todas las cifras que necesitan las preguntas de la gerencia.}

\textbf{¿Columna o medida?} Una \textbf{columna calculada} da un valor por
fila y se puede usar en un eje o en un segmentador. Una \textbf{medida} da un
solo número que se recalcula con los filtros de cada objeto visual.
\vez{Clase 13, ``Tres formas de calcular''}

\subtitulo{Columnas calculadas, en \texttt{Ventas}}

\begin{dax}
venta\_neta = Ventas[cantidad] * Ventas[precio\_unitario\_cop] - Ventas[descuento\_cop]\\
costo = Ventas[cantidad] * RELATED(Productos[costo\_unitario\_cop])\\
hora\_num = HOUR(Ventas[hora])\\
tipo\_cliente = IF(ISBLANK(Ventas[id\_cliente]), "Ocasional", "Socio")
\end{dax}

{\small \texttt{RELATED} trae a cada venta el costo de su producto: funciona
porque existe la relación \texttt{Productos}--\texttt{Ventas}. \texttt{IF}
con \texttt{ISBLANK} pregunta si la celda está vacía. \vez{Clase 13, paso 17;
Clase 14, paso 1}}

\subtitulo{Medidas, en \texttt{Ventas}}

Clic derecho en \texttt{Ventas} $>$ \menu{Nueva medida}{New measure}, una por
una. Da formato de porcentaje a las que terminan en \texttt{\%} y a
\texttt{Variación mensual} y \texttt{Cumplimiento}.

\begin{dax}
Ventas netas = SUM(Ventas[venta\_neta])\\
Costo total = SUM(Ventas[costo])\\
Margen = [Ventas netas] - [Costo total]\\
Margen \% = DIVIDE([Margen], [Ventas netas])\\
Unidades = SUM(Ventas[cantidad])\\
Tickets = DISTINCTCOUNT(Ventas[id\_ticket])\\
Ticket promedio = DIVIDE([Ventas netas], [Tickets])\\[4pt]
Ventas mes anterior = CALCULATE([Ventas netas], PREVIOUSMONTH(Calendario[Date]))\\
Variación mensual = IF(HASONEVALUE(Calendario[Año-Mes]),\\
\hspace*{2em}DIVIDE([Ventas netas] - [Ventas mes anterior], [Ventas mes anterior]))\\[4pt]
Meta de ventas = SUM(Tiendas[meta\_ventas\_mensual\_cop])\\
Cumplimiento = IF(HASONEVALUE(Calendario[Año-Mes]), DIVIDE([Ventas netas], [Meta de ventas]))\\
Meta de margen = 0.6\\
Ventas (con ceros) = COALESCE([Ventas netas], 0)
\end{dax}

\small
\renewcommand{\arraystretch}{1.15}
\begin{tabularx}{\textwidth}{@{}>{\raggedright\arraybackslash}p{3.6cm} >{\raggedright\arraybackslash}X >{\raggedright\arraybackslash}p{3.1cm}@{}}
\toprule
\textbf{Medida} & \textbf{Por qué se escribe así} & \textbf{Pista}\\
\midrule
\texttt{Tickets} & \texttt{DISTINCTCOUNT} cuenta compras, no líneas: un ticket con tres productos es una sola compra. & \vez{Clase 13, paso 18}\\
\texttt{Ticket promedio} & Cuánto deja, en promedio, cada compra. Responde a P3 y P6. & \\
\texttt{Variación mensual} & Solo tiene sentido con un único mes elegido; si no, queda en blanco. & \vez{Clase 15, paso 2}\\
\texttt{Meta de ventas} & La meta es \emph{mensual}: suma las metas de las tiendas visibles. Por eso \texttt{Cumplimiento} exige un único mes, igual que la variación. & \vez{Clase 15, paso 4}\\
\texttt{Ventas (con ceros)} & Un día sin ventas da vacío, y una línea de tiempo salta los vacíos sin avisar. Con ceros, el hueco se ve. & \vez{Clase 14, paso 11}\\
\bottomrule
\end{tabularx}
\renewcommand{\arraystretch}{1.0}
\normalsize

\clearpage
\begin{control}
\small Pon cada medida en una \menu{Tarjeta}{Card}, o todas en un objeto
visual \menu{Tabla}{Table}. \textbf{Sin ningún filtro} debes ver:

\vspace{2pt}
\renewcommand{\arraystretch}{1.15}
\begin{tabularx}{\linewidth}{@{}>{\raggedright\arraybackslash}X >{\raggedleft\arraybackslash}p{3.3cm} @{\hspace{1.2cm}} >{\raggedright\arraybackslash}X >{\raggedleft\arraybackslash}p{2.6cm}@{}}
\toprule
\textbf{Medida} & \textbf{Valor} & \textbf{Medida} & \textbf{Valor}\\
\midrule
Ventas netas & 766.082.130 & Unidades & 94.038\\
Costo total & 286.783.100 & Tickets & 41.489\\
Margen & 479.299.030 & Ticket promedio & 18.465\\
Margen \% & 62,6\% & Meta de ventas & 72.600.000\\
\bottomrule
\end{tabularx}
\renewcommand{\arraystretch}{1.0}

\vspace{4pt}
Con un objeto visual \menu{Tabla}{Table} que tenga \texttt{Calendario[Año-Mes]}
y las medidas, la fila de \textbf{2026-08} debe mostrar ventas netas de
\textbf{73.976.450} y margen de \textbf{63,3\%}; la de \textbf{2026-09},
\textbf{89.161.940} y \textbf{56,6\%}, con una variación mensual de
\textbf{20,5\%}.
\end{control}

\begin{cuidado}[Si tus ventas netas dan 766.255.930]
\small Te faltó quitar las 20 filas repetidas de \texttt{Ventas} (etapa 1):
esa diferencia, 173.800 pesos, son las ventas de un día que se cargó dos
veces.
\end{cuidado}

\begin{recuerda}[Puedes comparar con las medidas rápidas]
\small Si una fórmula no te sale, \menu{Nueva medida rápida}{New quick
measure} escribe una equivalente. Úsala para aprender, y luego compara su
DAX con el tuyo. \vez{Clase 15, paso 3}
\end{recuerda}

\begin{checklist}
  \item Las cuatro columnas calculadas están en \texttt{Ventas}.
  \item Las 13 medidas tienen el formato correcto (pesos sin decimales,
    porcentajes con uno).
  \item Las cifras de la caja verde coinciden con las mías.
  \item \texttt{Variación mensual} y \texttt{Cumplimiento} quedan en blanco
    cuando no hay un único mes elegido. Eso es lo correcto.
\end{checklist}

% ============================================================
\etapa{4}{Explorar: una página por cada pregunta}{encontrar, con evidencia, la respuesta a cada pregunta de la gerencia. Aquí todavía no importa que la página sea bonita.}

Crea una página llamada \textbf{Exploración}. Para cada pregunta, la tabla
te sugiere un objeto visual y los campos; \textbf{la respuesta la encuentras
tú}. Anota cada hallazgo en la ficha de la etapa 6 apenas lo veas.
\vez{Clase 14, pasos 4 a 16}

\small
\renewcommand{\arraystretch}{1.25}
\begin{tabularx}{\textwidth}{@{}>{\raggedright\arraybackslash}p{0.8cm} >{\raggedright\arraybackslash}p{4.6cm} >{\raggedright\arraybackslash}X@{}}
\toprule
& \textbf{Objeto visual y campos} & \textbf{Qué mirar}\\
\midrule
\textbf{P1} & Dos \textbf{líneas} separadas, por \texttt{Calendario[Inicio de mes]}: una con \texttt{Ventas netas}, otra con \texttt{Margen \%}. & ¿En qué mes se separan las dos curvas? Luego, una \textbf{matriz} con \texttt{Productos[categoria]} en filas, \texttt{Calendario[Año-Mes]} en columnas y \texttt{Margen \%} y \texttt{Unidades} en valores: compara agosto con septiembre. ¿Qué categoría cambió, y cuánto? ¿Y si quitas esa categoría con un filtro, cuánto crecieron las ventas?\\
\textbf{P2} & \textbf{Barras} de \texttt{Cumplimiento} por \texttt{Tiendas[tienda]}, con un \textbf{segmentador} de \texttt{Calendario[Año-Mes]}. & Recorre los meses de abril a septiembre. ¿Qué tiendas quedan por debajo del 100\% una y otra vez? Para Chapinero, unas \textbf{columnas} de \texttt{Ventas netas} por \texttt{Calendario[Día semana]}, filtradas a esa tienda, y compáralas con otra tienda.\\
\textbf{P3} & \textbf{Columnas apiladas al 100\%} de \texttt{Ventas netas} por \texttt{Calendario[Inicio de mes]} con \texttt{Ventas[canal]} en la leyenda. & ¿Cómo cambia la parte de la app desde marzo? Luego una \textbf{tabla} con \texttt{canal}, \texttt{Tickets} y \texttt{Ticket promedio}: ¿quién compra más por visita?\\
\textbf{P4} & \textbf{Columnas agrupadas} de \texttt{Unidades} por \texttt{Ventas[hora\_num]}, con \texttt{Productos[categoria]} en la leyenda. & ¿Cuáles son las horas de más movimiento? ¿Qué categoría manda en la mañana y cuál en la tarde?\\
\textbf{P5} & \textbf{Línea} de \texttt{Ventas netas} por \texttt{Calendario[Inicio de mes]}, filtrada a Usaquén, y la \texttt{Meta de ventas} como segunda línea. & ¿Desde qué mes supera su meta? Ojo: una tienda nueva no se compara, sin más, con una que lleva años abierta.\\
\textbf{P6} & \textbf{Tabla} con \texttt{Ventas[tipo\_cliente]}, \texttt{Tickets}, \texttt{Ticket promedio} y \texttt{Ventas netas}. & ¿Cuánto más deja un ticket de socio? ¿Qué parte de las ventas viene de socios? Y en \texttt{Clientes}: ¿cuántos socios aceptan promociones?\\
\bottomrule
\end{tabularx}
\renewcommand{\arraystretch}{1.0}
\normalsize

\subtitulo{El encargo extra: ¿qué pasó en Granada en julio?}

\begin{enumerate}[itemsep=2pt, topsep=2pt]
  \item Una \textbf{línea} con \texttt{Calendario[Date]} en el eje X y
    \texttt{Ventas (con ceros)} en el eje Y.
  \item Filtro de objeto visual: \texttt{Tiendas[tienda]} = Altavista
    Granada, y \texttt{Calendario[Año-Mes]} = 2026-07.
  \item Panel \menu{Análisis}{Analytics} $>$ \menu{Buscar
    anomalías}{Find anomalies}. \vez{Clase 14, paso 11}
\end{enumerate}

Cuando encuentres el problema, \textbf{valídalo antes de concluir}, como con
Envigado en la Clase 14: ¿pasa en las otras tiendas esos mismos días? ¿Hay
alguna venta, aunque sea una? ¿Cuántos días son? Con lo que muestran los
datos, ¿puedes saber si la tienda cerró o si la caja dejó de registrar? Si
no, tu insight es una \textbf{pregunta para la operación}, y así debes
presentarlo.

\begin{recuerda}[{{Un mapa, para la página final}}]
\small Antes de pasar a la etapa 5, prepara la ubicación. En la vista Tabla,
selecciona \texttt{Tiendas[latitud]} y en \menu{Herramientas de
columnas}{Column tools} $>$ \menu{Categoría de datos}{Data category} elige
\textbf{Latitud}; igual con \texttt{longitud} (\textbf{Longitud}) y
\texttt{ciudad} (\textbf{Ciudad}). Prueba un objeto visual \textbf{Azure
Maps} con \texttt{latitud}, \texttt{longitud} y \texttt{Ventas netas} como
tamaño de la burbuja. \vez{Clase 14, pasos 2 y 13}
\end{recuerda}

% ============================================================
\etapa{5}{La página para la gerencia}{construir una sola página que responda las preguntas sin que tengas que explicarla, y que la gerencia pueda filtrar por su cuenta.}

Crea una página nueva llamada \textbf{Comité de octubre}. Este es el boceto;
la ubicación exacta la decides tú, pero respeta el orden: arriba lo que
filtra y los indicadores, abajo el detalle. \vez{Clase 15, pasos 1 y 14}

\begin{center}
\begin{tikzpicture}[
  zona/.style={draw=esumerNavy!60, line width=0.6pt, rounded corners=2pt,
               fill=white, align=center, font=\scriptsize, inner sep=2pt},
  seg/.style={zona, fill=esumerTeal!12},
  ind/.style={zona, fill=esumerOrange!10},
]
  \draw[esumerNavy, line width=1pt, rounded corners=3pt, fill=esumerBg] (0,0) rectangle (16.8,9.2);
  \node[anchor=west, font=\footnotesize\bfseries, text=esumerNavy] at (0.3,8.75)
    {¿Cumplimos la meta de ventas, y con qué margen?};
  \node[seg, minimum width=2.9cm, minimum height=0.6cm] at (10.25,8.75) {Segmentador: Año-Mes};
  \node[seg, minimum width=2.3cm, minimum height=0.6cm] at (13.1,8.75) {Segm.: ciudad};
  \node[seg, minimum width=2.2cm, minimum height=0.6cm] at (15.5,8.75) {Segm.: canal};
  \node[ind, minimum width=3.9cm, minimum height=1.75cm] at (2.25,7.2) {\textbf{KPI}\\Ventas netas\\frente a Meta de ventas};
  \node[ind, minimum width=3.9cm, minimum height=1.75cm] at (6.35,7.2) {\textbf{KPI}\\Margen \%\\frente a Meta de margen};
  \node[ind, minimum width=3.9cm, minimum height=1.75cm] at (10.45,7.2) {\textbf{Tarjeta}\\Variación mensual};
  \node[ind, minimum width=3.9cm, minimum height=1.75cm] at (14.55,7.2) {\textbf{Tarjeta}\\Ticket promedio};
  \node[zona, minimum width=8.0cm, minimum height=2.9cm] at (4.3,4.55)
    {\textbf{Barras}: Cumplimiento por tienda\\(con una línea de referencia en 100\%)};
  \node[zona, minimum width=7.95cm, minimum height=2.9cm] at (12.5,4.55)
    {\textbf{Líneas}: Ventas netas por Inicio de mes\\(una línea por canal)};
  \node[zona, minimum width=8.0cm, minimum height=2.75cm] at (4.3,1.6)
    {\textbf{Azure Maps}: tiendas por coordenadas,\\burbuja = Ventas netas};
  \node[zona, minimum width=7.95cm, minimum height=2.75cm] at (12.5,1.6)
    {\textbf{Barras con jerarquía}: Margen \% por\\categoría $>$ subcategoría $>$ producto};
\end{tikzpicture}
\end{center}

\subtitulo{Lo que tienes que lograr}

\begin{enumerate}[itemsep=3pt, topsep=2pt]
  \item \textbf{Dos KPI.} \menu{Valor}{Value}, \menu{Eje de
    tendencia}{Trend axis} = \texttt{Calendario[Inicio de mes]} y
    \menu{Destino}{Target}. Los dos con dirección \menu{Alto es bueno}{High
    is good}. \hfill \vez{Clase 15, paso 5}
  \item \textbf{Tres segmentadores}: \texttt{Calendario[Año-Mes]} (lista
    desplegable, selección única), \texttt{Tiendas[ciudad]} y
    \texttt{Ventas[canal]}. \hfill \vez{Clase 15, pasos 7 y 8}
  \item \textbf{Una jerarquía} en \texttt{Productos}: \texttt{categoria}
    $>$ \texttt{subcategoria} $>$ \texttt{producto}, y explórala con las
    flechas del objeto visual. \hfill \vez{Clase 15, pasos 11 y 12}
  \item \textbf{Un filtro de objeto visual}: en las barras con
    jerarquía, usa \menu{N superior}{Top N}: los 5 valores de
    \texttt{Productos[producto]} con más
    \texttt{Ventas netas}, para que al bajar al nivel de producto no se
    llene la barra. \hfill \vez{Clase 15, paso 10}
  \item \textbf{Interacciones}, pensando qué compara cada objeto visual
    (\menu{Formato}{Format} $>$ \menu{Editar interacciones}{Edit
    interactions}). \vez{Clase 15, paso 13}
    \begin{itemize}[itemsep=1pt, topsep=1pt]
      \item \textbf{Mes}: deja \textbf{Filtrar} en los KPI (muestran el
        mes elegido) y pon \textbf{Ninguno} en la línea de ventas, para que
        siga mostrando todo el año.
      \item \textbf{Ciudad}: deja \textbf{Filtrar} en todo.
        \texttt{Meta de ventas} suma las metas de las tiendas visibles, así
        que la meta también se ajusta a la ciudad.
      \item \textbf{Canal}: la meta es de la tienda completa, no de un
        canal. Si filtras el KPI de ventas por canal, compararías solo la app
        con la meta de toda la tienda. Pon ahí \textbf{Ninguno}.
    \end{itemize}
\end{enumerate}

\begin{control}
\small Sin segmentadores, el \textbf{KPI de ventas} muestra el último mes,
septiembre: \textbf{89.161.940} frente a la meta de \textbf{72.600.000}. El
\textbf{KPI de margen} muestra \textbf{56,6\%} frente a \textbf{60\%}.
Con el segmentador en \textbf{2026-09}, la tarjeta de variación mensual
dice \textbf{20,5\%} y las barras de cumplimiento van de \textbf{103,3\%}
(Chapinero) a \textbf{163,2\%} (Usaquén).
\end{control}

\begin{cuidado}[{{Ventas en verde, margen en rojo: no es un error}}]
\small Los dos KPI de septiembre no se contradicen: responden preguntas
distintas. Si uno está verde y el otro rojo, ahí empieza tu respuesta a P1.
Antes de presentar, prueba la página como alguien que no la conoce: elige
un mes, una ciudad y un canal, y comprueba que cada objeto visual cambia
como esperas. \vez{Clase 15, paso 15}
\end{cuidado}

\subtitulo{La página está lista cuando\ldots}

\begin{checklist}
  \item El título dice qué pregunta responde la página, y cada objeto
    visual tiene un título que se entiende sin ti.
  \item Los dos KPI tienen meta y dirección \textbf{Alto es bueno} y, sin
    un mes elegido, muestran la tendencia de los meses.
  \item Al elegir \textbf{Bogotá} en el segmentador de ciudad, la meta del KPI
    de ventas baja a \textbf{24.000.000} (Chapinero más Usaquén).
  \item Al elegir \textbf{App} en el segmentador de canal, el KPI de ventas
    no cambia, y los demás objetos visuales sí.
  \item La jerarquía baja de categoría a subcategoría y a producto, y vuelve
    a subir.
  \item El mapa muestra seis burbujas en cuatro ciudades.
  \item Los colores significan lo mismo en toda la página: si \texttt{App}
    es naranja en la línea, lo es en todas partes (Clase 14).
  \item Cada objeto visual responde a una de las preguntas de la gerencia.
    Si alguno no responde a ninguna, sobra o va en la página de exploración.
\end{checklist}

% ============================================================
\etapa{6}{Comunicar: ficha de insights y presentación}{convertir lo que encontraste en decisiones, y contarlo en cinco minutos al comité.}

\textbf{Un insight no es una cifra.} Es una frase con evidencia, que alguien
puede usar para decidir. Antes de escribirlo, pásalo por las cinco preguntas
de validación de la Clase 14: ¿los totales cuadran?, ¿el periodo está
completo?, ¿la comparación es justa?, ¿hay suficientes datos?, ¿es causa o
solo relación? \vez{Clase 14, Hora 3}

\begin{center}
\begin{tikzpicture}[
  pieza/.style={caja, minimum width=3.8cm, text width=3.5cm, minimum height=2.0cm},
  tit/.style={font=\bfseries\small, text=white, fill=esumerTeal, rounded corners=2pt,
              minimum width=3.8cm, inner sep=3pt, text height=1.6ex, text depth=0.4ex},
]
  \foreach \x/\tit in {2.2/Pregunta, 6.5/Evidencia, 10.8/Insight, 15.1/Recomendación}
    {\node[tit] at (\x,1.4) {\tit};}
  \node[pieza] (a) at (2.2,0) {\small ¿A qué horas reforzamos el personal?};
  \node[pieza] (b) at (6.5,0) {\small Columnas de unidades por hora y categoría};
  \node[pieza] (c) at (10.8,0) {\small La frase con tu cifra: qué pasa, cuánto y cuándo};
  \node[pieza] (d) at (15.1,0) {\small Qué hacer, quién lo hace y cómo saber si funcionó};
  \draw[flecha] (a) -- (b);
  \draw[flecha] (b) -- (c);
  \draw[flecha] (c) -- (d);
\end{tikzpicture}
\end{center}

\subtitulo{Ficha de insights}

\renewcommand{\arraystretch}{3.0}
\begin{tabularx}{\textwidth}{@{}>{\centering\arraybackslash}p{0.7cm} >{\raggedright\arraybackslash}X >{\raggedright\arraybackslash}X >{\raggedright\arraybackslash}X >{\centering\arraybackslash}p{1.5cm}@{}}
\toprule
\textbf{P} & \textbf{Evidencia (cifra y objeto visual)} & \textbf{Insight} & \textbf{Recomendación} & \textbf{Validado}\\
\midrule
& & & & \\ \hline
& & & & \\ \hline
& & & & \\ \hline
& & & & \\ \hline
& & & & \\
\bottomrule
\end{tabularx}\par
\renewcommand{\arraystretch}{1.0}

{\footnotesize\itshape En la columna P, escribe el número de la pregunta de
la gerencia, o \textbf{D} si es un hallazgo sobre la calidad de los datos
(las filas repetidas, Granada).}

\subtitulo{La presentación: cinco minutos para el comité}

\begin{center}
\begin{tikzpicture}[
  bloque/.style={draw=none, rounded corners=3pt, text=white, align=center,
                 minimum height=1.6cm, inner sep=3pt, font=\small},
]
  \node[bloque, fill=esumerNavy, minimum width=2.6cm, text width=2.4cm] (b1) at (1.3,0)
    {\textbf{30 s}\\[1pt]{\scriptsize El contexto y las preguntas}};
  \node[bloque, fill=esumerSky, minimum width=7.6cm, text width=7.3cm] (b2) at (6.55,0)
    {\textbf{3 min}\\[1pt]{\scriptsize Tres insights, cada uno con su evidencia en la
     página de Power BI. Empieza por el que lleva a la decisión más importante}};
  \node[bloque, fill=esumerOrange, minimum width=3.4cm, text width=3.2cm] (b3) at (12.15,0)
    {\textbf{1 min}\\[1pt]{\scriptsize Recomendaciones: qué hacer el próximo mes}};
  \node[bloque, fill=esumerMagenta, minimum width=2.8cm, text width=2.6cm] (b4) at (15.4,0)
    {\textbf{30 s}\\[1pt]{\scriptsize Lo que los datos no dicen y qué falta saber}};
\end{tikzpicture}
\end{center}

\small
\begin{itemize}[itemsep=1pt, topsep=2pt]
  \item Presenta \textbf{desde la página de Power BI}, filtrando en vivo,
    no con capturas de pantalla.
  \item El \textbf{título} de la página dice qué pregunta responde; la
    \textbf{conclusión} la dices tú, con cifras (Clase 15, paso 1).
  \item Un compañero hace de gerente: puede interrumpirte con una pregunta
    que no estaba en la carta. Respóndela filtrando, o di qué dato te falta.
\end{itemize}
\normalsize

% ============================================================
\clearpage
\section*{\color{esumerNavy}Autoevaluación: ¿tu tablero está listo?}

Antes de entregar, revisa cada criterio con honestidad y marca en la
última columna dónde estás.

\vspace{2pt}
\small
\renewcommand{\arraystretch}{1.3}
\begin{tabularx}{\textwidth}{@{}>{\raggedright\arraybackslash}p{2.6cm} >{\raggedright\arraybackslash}X >{\raggedright\arraybackslash}X >{\centering\arraybackslash}p{0.9cm}@{}}
\toprule
\textbf{Criterio} & \textbf{En proceso} & \textbf{Logrado} & \textbf{Yo}\\
\midrule
\textbf{Datos} & Hay pasos de limpieza, pero quedan variantes, duplicados o tipos equivocados. & Las cuatro tablas coinciden con las cifras de control de la etapa 1. & \\
\textbf{Modelo} & Faltan relaciones o alguna es muchos a muchos; no hay \texttt{Calendario}. & Modelo estrella con cuatro relaciones *:1 de dirección única, \texttt{Calendario} marcada. & \\
\textbf{Medidas} & Las cifras se calculan con columnas o resúmenes automáticos, o no cuadran. & Las 13 medidas existen, tienen formato y coinciden con la etapa 3. & \\
\textbf{Exploración} & Hay gráficos, pero no responden a una pregunta concreta. & Cada pregunta de la gerencia tiene un objeto visual que la responde. & \\
\textbf{Página final} & Los filtros no funcionan como se espera, o la página está saturada. & KPI con meta, segmentadores, jerarquía, mapa e interacciones revisadas. & \\
\textbf{Insights} & Describen gráficos (``las ventas suben en diciembre''). & Cada uno tiene cifra, comparación, explicación y recomendación, y fue validado. & \\
\textbf{Calidad} & Las filas repetidas y Granada no se mencionan. & Se identifican, se explica qué se hizo y qué queda por preguntar. & \\
\textbf{Presentación} & Más de cinco minutos, o se lee la pantalla. & Cinco minutos, desde Power BI, terminando en decisiones. & \\
\bottomrule
\end{tabularx}
\renewcommand{\arraystretch}{1.0}
\normalsize

\subtitulo{Si algo no cuadra}

\small
\renewcommand{\arraystretch}{1.2}
\begin{tabularx}{\textwidth}{@{}>{\raggedright\arraybackslash}p{5.6cm} >{\raggedright\arraybackslash}X@{}}
\toprule
\textbf{Lo que ves} & \textbf{Causa más probable}\\
\midrule
La misma cifra en todas las filas de una tabla o en todas las barras. & Falta la relación entre esa dimensión y \texttt{Ventas}, o va en la dirección contraria (etapa 2).\\
\texttt{categoria} sigue con más de 4 valores. & Faltó \menu{Recortar}{Trim}, o el reemplazo se hizo sin pasar antes a MAYÚSCULAS.\\
Ventas netas de 766.255.930. & No quitaste los duplicados de \texttt{Ventas}.\\
\texttt{Cumplimiento} o \texttt{Variación mensual} en blanco. & No hay un único mes elegido. Es lo esperado: usa el segmentador de mes.\\
Al elegir un canal, el KPI de ventas queda muy lejos de la meta. & La meta es de toda la tienda, no de un canal: pon en \textbf{Ninguno} la interacción del segmentador de canal sobre ese KPI.\\
\texttt{costo} da error. & Falta la relación \texttt{Productos}--\texttt{Ventas}: \texttt{RELATED} la necesita.\\
\texttt{hora\_num} da error. & \texttt{hora} quedó como texto. Cámbiala a tipo \textbf{Hora} en Power Query.\\
El mapa no muestra nada. & Falta la categoría de datos de \texttt{latitud} y \texttt{longitud}, o el equipo no tiene internet o tiene Azure Maps deshabilitado.\\
\bottomrule
\end{tabularx}
\renewcommand{\arraystretch}{1.0}
\normalsize

\subtitulo{Documentación oficial}

\small
\begin{itemize}[itemsep=1pt, topsep=2pt]
  \item Etapa 1: \href{\urlconsultas}{\color{esumerOrange}Información general de consultas en Power BI Desktop}.
  \item Etapa 2: \href{\urlestrella}{\color{esumerOrange}Comprender el esquema de estrella} y
    \href{\urlrelaciones}{\color{esumerOrange}Crear y administrar relaciones}.
  \item Etapa 3: \href{\urlcolumnas}{\color{esumerOrange}Columnas calculadas} y
    \href{\urlmedidas}{\color{esumerOrange}Medidas en Power BI Desktop}.
  \item Etapas 4 y 5: \href{\urlkpi}{\color{esumerOrange}Objeto visual KPI},
    \href{\urlsegmentadores}{\color{esumerOrange}Segmentadores},
    \href{\urlfiltros}{\color{esumerOrange}Filtros},
    \href{\urlinteracciones}{\color{esumerOrange}Interacciones},
    \href{\urlexplorar}{\color{esumerOrange}Explorar en profundidad} y
    \href{\urlazuremaps}{\color{esumerOrange}Azure Maps}.
  \item Diseño de la página: \href{\urltrucos}{\color{esumerOrange}Sugerencias y trucos para crear informes}.
\end{itemize}
\normalsize

\end{document}
